\chapter{Categories, Functors, and Presheaves: A Primer}\label{ch02:chap}

\section{Purpose and scope}\label{ch02:sec:purpose}

This chapter develops, from nothing, the standard category theory used in
the rest of the monograph: categories, functors, natural transformations,
opposite categories, presheaves, representable presheaves, the Yoneda
lemma, and limits with particular attention to pullbacks. Everything in
this chapter is textbook material
\cite{awodey,riehl,maclanecwm,maclane}, presented as such; none of it is a
project contribution, and the chapter's only originality is in the choice
and arrangement of examples. Readers fluent in the subject should skim for
notation --- in particular the poset conventions of
Section~\ref{ch02:sec:posets-as-cats} and the small posets introduced in
Section~\ref{ch02:sec:laboratory}, which recur throughout Part~II --- and
move on.

The pedagogical strategy is dictated by the destination. The finite theory
of Part~II lives entirely on a sixteen-element powerset lattice, viewed as
a category; every presheaf that matters there is a functor out of a finite
poset, every pullback that matters is an intersection of subsets, and
every application of the Yoneda lemma is an application of its poset case.
This chapter therefore works the poset case of everything in full, on
posets small enough that all objects can be enumerated by hand: the
two-element chain, the three-element chain, a three-element ``vee,'' and
the four-element powerset lattice of a two-element set. The general
definitions are always given --- the legacy category $\cC$ of Chapter~5 is
not a poset, and the countermodel of Chapter~6 lives in a category of
finite sets and injections --- but the general theory is cited rather than
redeveloped where the poset case suffices for Part~II. A reader wanting a
refresher on elementary proof technique before starting may consult
Hammack \cite{bookofproof}; the proofs below use nothing beyond induction,
case analysis, and careful bookkeeping with sets and functions.

\section{Categories}\label{ch02:sec:categories}

\begin{definition}[Category]\label{ch02:def:category}
A \emph{category} $\cC$ consists of the following data:
\begin{enumerate}
\item a class $\Ob(\cC)$ of \emph{objects};
\item for each ordered pair of objects $(A, B)$, a set
$\Hom_{\cC}(A, B)$ of \emph{morphisms} from $A$ to $B$ (written
$f : A \to B$; $A$ is the \emph{domain} of $f$ and $B$ its
\emph{codomain});
\item for each triple of objects $(A, B, C)$, a \emph{composition}
operation
\[
\circ : \Hom_{\cC}(B, C) \times \Hom_{\cC}(A, B) \to \Hom_{\cC}(A, C),
\qquad (g, f) \mapsto g \circ f ;
\]
\item for each object $A$, an \emph{identity} morphism
$\id_A \in \Hom_{\cC}(A, A)$;
\end{enumerate}
subject to two axioms:
\begin{enumerate}
\item[(C1)] \emph{Associativity}: $h \circ (g \circ f) = (h \circ g)
\circ f$ whenever the composites are defined;
\item[(C2)] \emph{Unitality}: $\id_B \circ f = f = f \circ \id_A$ for
every $f : A \to B$.
\end{enumerate}
\end{definition}

We assume morphism sets with distinct domain-codomain pairs are disjoint,
so that every morphism has a well-defined domain and codomain. A category
is \emph{small} when $\Ob(\cC)$ is a set (not a proper class) and
\emph{locally small} when each $\Hom_{\cC}(A,B)$ is a set, as
Definition~\ref{ch02:def:category} already requires. Size issues play one
role in this monograph: topos-existence statements require the site to be
small in a chosen universe, and the finite lattice $L$ of Part~II is not
merely small but finite, which is why exhaustive checking is possible at
all.

\begin{example}[The category of sets]\label{ch02:ex:set}
$\Set$ has all sets (in a fixed universe) as objects and all functions as
morphisms, with ordinary composition of functions and the ordinary
identity functions. Axioms (C1) and (C2) are the familiar facts that
composition of functions is associative and that identity functions are
neutral. $\Set$ is locally small but not small.
\end{example}

\begin{example}[Finite sets]\label{ch02:ex:finset}
$\mathbf{FinSet}$ has finite sets as objects and all functions between
them as morphisms; it is a \emph{full subcategory} of $\Set$, meaning its
morphism sets are exactly those of $\Set$ between the retained objects.
Variants with restricted morphisms also matter later: the category of
finite sets and \emph{injections} (same objects, only injective functions)
is a category, because identities are injective and composites of
injections are injections. Chapter~6's countermodel lives in the full
subcategory of that category on sets of size at most two. Note that
restricting morphisms can change categorical properties drastically ---
which is the point of the countermodel.
\end{example}

\subsection{Posets as categories}\label{ch02:sec:posets-as-cats}

\begin{definition}[Poset]\label{ch02:def:poset}
A \emph{partially ordered set} (poset) is a set $P$ with a relation
$\le$ that is reflexive ($U \le U$), transitive ($U \le V$ and $V \le W$
imply $U \le W$), and antisymmetric ($U \le V$ and $V \le U$ imply
$U = V$).
\end{definition}

\begin{proposition}[A poset is a category]\label{ch02:prop:poset-category}
Let $(P, \le)$ be a poset. Define a category, also written $P$, with
$\Ob(P) = P$ and
\[
\Hom_P(U, V) =
\begin{cases}
\{\, \iota_{U,V} \,\} & \text{if } U \le V, \\
\varnothing & \text{otherwise,}
\end{cases}
\]
where $\iota_{U,V}$ is a formal symbol (``the'' morphism from $U$ to
$V$). Composition is $\iota_{V,W} \circ \iota_{U,V} = \iota_{U,W}$ and
the identity on $U$ is $\iota_{U,U}$. Then $P$ is a category.
\end{proposition}

\begin{proof}
Composition is well defined: if $U \le V$ and $V \le W$ then $U \le W$ by
transitivity, so $\iota_{U,W}$ exists. Identities exist by reflexivity.
Associativity (C1): both $h \circ (g \circ f)$ and $(h \circ g) \circ f$,
when defined, are morphisms $U \to Z$ for the same pair $(U, Z)$, and
$\Hom_P(U, Z)$ has at most one element, so they are equal. Unitality (C2)
holds for the same reason: $\id_V \circ \iota_{U,V}$ and $\iota_{U,V}$
are both elements of the at-most-one-element set $\Hom_P(U,V)$.
\end{proof}

We isolate the at-most-one-morphism property just used, because it
is invoked silently on almost every page of Part~II.

\begin{proposition}[All diagrams in a poset commute]\label{ch02:prop:poset-commute}
In a poset regarded as a category, any two morphisms with the same domain
and the same codomain are equal. Consequently every diagram of morphisms
commutes: any two composites with the same source and target agree.
\end{proposition}

\begin{proof}
Immediate from $|\Hom_P(U,V)| \le 1$: parallel morphisms lie in a common
set with at most one element. A composite of morphisms $U \to \cdots \to
V$ is a morphism $U \to V$, so any two such composites are parallel, hence
equal.
\end{proof}

\begin{example}[The standard small posets]\label{ch02:ex:small-posets}
Four posets recur throughout this monograph.
\begin{enumerate}
\item The \emph{two-element chain} $[2] = \{0, 1\}$ with $0 \le 1$. As a
category it has two objects and three morphisms: $\id_0$, $\id_1$, and
$f : 0 \to 1$.
\item The \emph{three-element chain} $[3] = \{0, 1, 2\}$ with
$0 \le 1 \le 2$. Six morphisms: three identities and the arrows
$0 \to 1$, $1 \to 2$, $0 \to 2$, with $(1 \to 2) \circ (0 \to 1) =
(0 \to 2)$ --- the only nontrivial composition, forced by
Proposition~\ref{ch02:prop:poset-commute}.
\item The \emph{vee} $\Lambda = \{x, y, t\}$ with $x \le t$, $y \le t$,
and $x, y$ incomparable. Five morphisms: three identities, $x \to t$,
$y \to t$. There is no morphism between $x$ and $y$ in either direction.
\item The \emph{square} $B = \mathcal{P}(\{1, 2\})$, the powerset of a
two-element set ordered by inclusion: objects $\varnothing$, $\{1\}$,
$\{2\}$, $\{1,2\}$. Counting ordered pairs $U \subseteq V$:
$\varnothing$ is below all four objects, $\{1\}$ below itself and
$\{1,2\}$, $\{2\}$ likewise, and $\{1,2\}$ below itself only ---
$4 + 2 + 2 + 1 = 9$ morphisms.
\end{enumerate}
\end{example}

\begin{proposition}[Counting inclusions in a powerset]\label{ch02:prop:threepow}
For a finite set $X$, the number of ordered pairs $(U, V)$ with
$U \subseteq V \subseteq X$ --- that is, the number of morphisms of the
category $\mathcal{P}(X)$ --- is $3^{|X|}$.
\end{proposition}

\begin{proof}
A pair $U \subseteq V$ is determined by assigning to each element of $X$
one of three mutually exclusive statuses: in $U$ (hence in $V$), in
$V$ but not $U$, or outside $V$. Distinct assignments give distinct pairs
and every pair arises, so the count is $3^{|X|}$.
\end{proof}

For $|X| = 2$ this gives the $9$ morphisms of $B$ just counted by hand;
for $|X| = 4$ it gives $3^4 = 81$, which is the arrow count of the context
lattice $L = \mathcal{P}(\{0,1\} \times \{0,1\})$ recorded by the finite
certificate and used constantly in Part~II. The agreement between this
one-line combinatorial argument and the certificate's exhaustive
enumeration is a small instance of the corroboration discipline of
Chapter~11.

\subsection{Monoids as one-object categories}\label{ch02:sec:monoids}

\begin{definition}[Monoid]\label{ch02:def:monoid}
A \emph{monoid} is a set $M$ with an associative binary operation
$(m, m') \mapsto m m'$ and a two-sided unit $e \in M$: $e m = m = m e$.
\end{definition}

\begin{proposition}[Monoids are one-object categories]\label{ch02:prop:monoid-category}
For a monoid $M$, define a category $\mathbf{B}M$ with a single object
$\ast$, $\Hom(\ast, \ast) = M$, composition $m' \circ m = m' m$
(multiplication), and $\id_\ast = e$. Then $\mathbf{B}M$ is a category.
Conversely, if $\cC$ is a category with exactly one object $\ast$, then
$\Hom_{\cC}(\ast, \ast)$ with composition as multiplication and
$\id_\ast$ as unit is a monoid, and the two constructions are mutually
inverse.
\end{proposition}

\begin{proof}
For $\mathbf{B}M$: all composites are defined (there is one object), and
(C1), (C2) are exactly associativity and unitality of $M$. Conversely, in
a one-object category the composition operation is a total binary
operation on the single morphism set; (C1) makes it associative and (C2)
makes $\id_\ast$ a two-sided unit. Applying the first construction to the
monoid so obtained returns the original category, and applying the second
to $\mathbf{B}M$ returns $M$.
\end{proof}

Posets and monoids are the two degenerate extremes of the notion of
category: a poset has many objects but at most one morphism between any
two, so all information is in \emph{which} morphisms exist; a monoid has
one object and many morphisms, so all information is in \emph{how}
morphisms compose. General categories mix both kinds of information. The
finite theory of Part~II lives at the poset extreme; the legacy category
$\cC$ of Chapter~5 does not, which is why the general definitions are
needed at all.

\subsection{Isomorphisms, monomorphisms, epimorphisms}\label{ch02:sec:special-morphisms}

\begin{definition}\label{ch02:def:iso-mono-epi}
Let $f : A \to B$ be a morphism in a category $\cC$.
\begin{enumerate}
\item $f$ is an \emph{isomorphism} if there exists $g : B \to A$ with
$g \circ f = \id_A$ and $f \circ g = \id_B$; such $g$ is unique when it
exists and is written $f^{-1}$.
\item $f$ is a \emph{monomorphism} (mono) if for all parallel
$h, k : Z \to A$, $f \circ h = f \circ k$ implies $h = k$.
\item $f$ is an \emph{epimorphism} (epi) if for all parallel
$h, k : B \to Z$, $h \circ f = k \circ f$ implies $h = k$.
\end{enumerate}
\end{definition}

\begin{proposition}[Monos and epis in $\Set$]\label{ch02:prop:set-mono-epi}
In $\Set$, a function is a monomorphism if and only if it is injective,
and an epimorphism if and only if it is surjective.
\end{proposition}

\begin{proof}
\emph{Mono $\Leftrightarrow$ injective.} If $f : A \to B$ is injective
and $f \circ h = f \circ k$ for $h, k : Z \to A$, then for each
$z \in Z$, $f(h(z)) = f(k(z))$, so $h(z) = k(z)$; thus $h = k$.
Conversely, suppose $f$ is a mono and $f(a) = f(a')$. Let
$Z = \{\ast\}$ be a one-element set and define $h, k : Z \to A$ by
$h(\ast) = a$, $k(\ast) = a'$. Then $f \circ h = f \circ k$, so
$h = k$, so $a = a'$.

\emph{Epi $\Leftrightarrow$ surjective.} If $f$ is surjective and
$h \circ f = k \circ f$, then every $b \in B$ is $f(a)$ for some $a$, so
$h(b) = h(f(a)) = k(f(a)) = k(b)$; thus $h = k$. Conversely, suppose $f$
is not surjective; choose $b_0 \in B$ with $b_0 \notin \im f$. Define
$h, k : B \to \{0, 1\}$ by $h$ constantly $0$, and $k(b) = 1$ if
$b = b_0$, else $k(b) = 0$. Then $h \ne k$ but $h \circ f = k \circ f$,
since $f$ never hits $b_0$. So $f$ is not an epi.
\end{proof}

\begin{proposition}[Isomorphisms in a poset are identities]\label{ch02:prop:poset-iso}
In a poset regarded as a category, the only isomorphisms are the
identity morphisms.
\end{proposition}

\begin{proof}
Suppose $f : U \to V$ has an inverse $g : V \to U$. The existence of $f$
and $g$ means $U \le V$ and $V \le U$, so $U = V$ by antisymmetry. Then
$f \in \Hom_P(U, U) = \{\id_U\}$, so $f = \id_U$.
\end{proof}

This innocuous fact has a consequence used in
Section~\ref{ch02:sec:yoneda}: in a poset, isomorphic representables have
equal representing objects --- there is no ``isomorphism slack.''

\subsection{Split epimorphisms and section-retraction pairs}\label{ch02:sec:split}

The next definition is elementary, but it is the single most consequential
definition of the chapter for this monograph's story: the legacy collapse
of Chapter~5 turns on it.

\begin{definition}[Section, retraction, split epi]\label{ch02:def:split}
A morphism $\varepsilon : B \to A$ is a \emph{split epimorphism} if there
exists $\eta : A \to B$ with $\varepsilon \circ \eta = \id_A$. One then
calls $\eta$ a \emph{section} of $\varepsilon$, calls $\varepsilon$ a
\emph{retraction} of $\eta$, and calls $(\eta, \varepsilon)$ a
\emph{section-retraction pair}. Nothing is assumed about the other
composite $\eta \circ \varepsilon : B \to B$, which in general is not
$\id_B$.
\end{definition}

\begin{proposition}\label{ch02:prop:split-epi-is-epi}
Every split epimorphism is an epimorphism, and every section is a
monomorphism.
\end{proposition}

\begin{proof}
Let $\varepsilon \circ \eta = \id_A$. If
$h \circ \varepsilon = k \circ \varepsilon$, compose on the right with
$\eta$: $h = h \circ \id_A = h \circ \varepsilon \circ \eta =
k \circ \varepsilon \circ \eta = k$. Dually, if
$\eta \circ h = \eta \circ k$, compose on the left with $\varepsilon$ to
get $h = k$.
\end{proof}

\begin{example}\label{ch02:ex:split-in-set}
In $\Set$, every surjection $\varepsilon : B \to A$ with $A$ nonempty (or
$B = A = \varnothing$) is split: choose, for each $a \in A$, an element
$\eta(a) \in \varepsilon^{-1}(a)$. The composite $\eta \circ \varepsilon$
sends each $b$ to the chosen representative of its fiber --- generally
not $b$ itself. This is the shape to keep in mind: passing from $B$ down
to $A$ and back up loses exactly the within-fiber information. The
registered RPM data of the legacy construction have this shape:
$\eta_X : X \to R(X)$, $\varepsilon_X : R(X) \to X$,
$\varepsilon_X \circ \eta_X = \id_X$, with the composite the other way
explicitly \emph{not} required to be the identity --- a
section-retraction pair, not an adjoint equivalence, encoding the
program-wide rule that prerequisite satisfaction is one-directional.
\end{example}

\begin{remark}\label{ch02:rem:no-split-in-posets}
In a poset, if $\varepsilon : V \to U$ is a split epimorphism with
section $\eta : U \to V$, then $U \le V$ and $V \le U$, so $U = V$ and
both maps are $\id_U$ (Proposition~\ref{ch02:prop:poset-iso}). Posets
harbor no nontrivial section-retraction pairs --- one reason the
successor's context factor, a poset, is structurally immune to the
split-epi collapse mechanism of Chapter~5.
\end{remark}

\begin{exercise}\label{ch02:exr:poset-monoepi}
Show that in a poset regarded as a category, \emph{every} morphism is
both a monomorphism and an epimorphism. (Hint:
Proposition~\ref{ch02:prop:poset-commute} makes the required implications
vacuous.) Conclude that ``mono and epi'' does not imply ``iso'' in
general categories.
\end{exercise}

\section{Functors}\label{ch02:sec:functors}

\begin{definition}[Functor]\label{ch02:def:functor}
Let $\cC$ and $\mathcal{D}$ be categories. A \emph{functor}
$F : \cC \to \mathcal{D}$ assigns to each object $A \in \Ob(\cC)$ an
object $F(A) \in \Ob(\mathcal{D})$ and to each morphism $f : A \to B$ a
morphism $F(f) : F(A) \to F(B)$, such that
\begin{enumerate}
\item[(F1)] $F(g \circ f) = F(g) \circ F(f)$ whenever $g \circ f$ is
defined, and
\item[(F2)] $F(\id_A) = \id_{F(A)}$ for every object $A$.
\end{enumerate}
\end{definition}

Functors compose in the evident way (apply one, then the other), and the
identity assignment is a functor; small categories and functors between
them themselves form a category \cite{maclanecwm}.

\begin{example}\label{ch02:ex:functors-basic}
(i) The inclusion $\mathbf{FinSet} \to \Set$ is a functor. (ii) For a
fixed set $S$, the assignment $A \mapsto A \times S$,
$f \mapsto f \times \id_S$ is a functor $\Set \to \Set$. (iii) For any
object $C$ of a locally small category $\cC$, the \emph{covariant hom
functor} $\Hom_{\cC}(C, -) : \cC \to \Set$ sends $B$ to
$\Hom_{\cC}(C, B)$ and $g : B \to B'$ to the map ``postcompose with
$g$.'' Functoriality is associativity and unitality of composition in
$\cC$.
\end{example}

\begin{proposition}[Functors between posets are monotone maps]\label{ch02:prop:functor-monotone}
Let $P$ and $Q$ be posets regarded as categories. A functor
$F : P \to Q$ is exactly the same data as a monotone (order-preserving)
function $F : P \to Q$: a function with $U \le V \Rightarrow
F(U) \le F(V)$.
\end{proposition}

\begin{proof}
Given a functor $F$, its object part is a function $P \to Q$. If
$U \le V$, the morphism $\iota_{U,V}$ exists, and $F(\iota_{U,V})$ is a
morphism $F(U) \to F(V)$ in $Q$, which by definition of $Q$ as a category
means $F(U) \le F(V)$. So the object part is monotone. Conversely, a
monotone function $F$ extends to morphisms in exactly one way (send
$\iota_{U,V}$ to $\iota_{F(U), F(V)}$, which exists by monotonicity), and
(F1), (F2) hold automatically by
Proposition~\ref{ch02:prop:poset-commute} applied in $Q$: all the
required equations are between parallel morphisms in a poset. The two
constructions are mutually inverse.
\end{proof}

\begin{example}\label{ch02:ex:functors-two-chain}
The functors $[2] \to [2]$ are the monotone self-maps of the two-element
chain: the constant map at $0$, the constant map at $1$, and the
identity. The non-monotone swap $0 \mapsto 1, 1 \mapsto 0$ is not a
functor: it would have to send the morphism $0 \to 1$ to a morphism
$1 \to 0$, and there is none.
\end{example}

\begin{proposition}[Functors between monoids are homomorphisms]\label{ch02:prop:functor-monoid}
For monoids $M$ and $N$, functors $\mathbf{B}M \to \mathbf{B}N$
correspond exactly to monoid homomorphisms $M \to N$ (functions
preserving multiplication and unit).
\end{proposition}

\begin{proof}
A functor sends $\ast$ to $\ast$ (no choice) and each morphism
$m \in M$ to a morphism $F(m) \in N$; (F1) says
$F(m' m) = F(m') F(m)$ and (F2) says $F(e_M) = e_N$ --- precisely the
homomorphism conditions. Conversely any homomorphism defines a functor
this way.
\end{proof}

\begin{definition}[Faithful, full]\label{ch02:def:faithful-full}
A functor $F : \cC \to \mathcal{D}$ is \emph{faithful} if each map
$\Hom_{\cC}(A, B) \to \Hom_{\mathcal{D}}(F A, F B)$ is injective,
\emph{full} if each is surjective, and \emph{fully faithful} if both.
\end{definition}

\section{Natural transformations}\label{ch02:sec:nat}

Functors compare categories; natural transformations compare functors.
The definition is short, but its content is the commuting square, and the
examples below are chosen to show that the square has real force.

\begin{definition}[Natural transformation]\label{ch02:def:nat}
Let $F, G : \cC \to \mathcal{D}$ be functors. A \emph{natural
transformation} $\alpha : F \Rightarrow G$ assigns to each object
$A \in \Ob(\cC)$ a morphism $\alpha_A : F(A) \to G(A)$ in $\mathcal{D}$
(the \emph{component} at $A$) such that for every morphism
$f : A \to B$ in $\cC$,
\[
G(f) \circ \alpha_A = \alpha_B \circ F(f)
\qquad \text{(naturality at } f\text{)}.
\]
If every component is an isomorphism, $\alpha$ is a \emph{natural
isomorphism}, and $F$ and $G$ are \emph{naturally isomorphic}, written
$F \cong G$.
\end{definition}

Natural transformations compose componentwise
($(\beta \circ \alpha)_A = \beta_A \circ \alpha_A$; naturality of the
composite follows by pasting the two squares), and the identity
transformation has components $\id_{F(A)}$. Hence:

\begin{definition}[Functor category]\label{ch02:def:functor-category}
For categories $\cC$ (small) and $\mathcal{D}$, the \emph{functor
category} $[\cC, \mathcal{D}]$ has functors $\cC \to \mathcal{D}$ as
objects and natural transformations as morphisms. Its isomorphisms are
exactly the natural isomorphisms.
\end{definition}

\begin{example}[Natural transformations into a poset are pointwise
inequalities]\label{ch02:ex:nat-poset}
Let $\cC$ be any small category and $Q$ a poset. For functors
$F, G : \cC \to Q$, a natural transformation $F \Rightarrow G$ exists if
and only if $F(A) \le G(A)$ for every object $A$, and is then unique.
Indeed, a component $\alpha_A$ is a morphism $F(A) \to G(A)$ in $Q$,
which exists iff $F(A) \le G(A)$ and is then unique; the naturality
squares hold automatically by
Proposition~\ref{ch02:prop:poset-commute} applied in $Q$. So the functor
category $[\cC, Q]$ is itself a poset: functors ordered pointwise. This
is the first hint of a theme that recurs with presheaves: when the target
or source collapses to a poset, natural-transformation data collapse to
order data.
\end{example}

\begin{example}[Intertwiners between monoid homomorphisms]\label{ch02:ex:nat-monoid}
Let $f, g : M \to N$ be monoid homomorphisms, regarded as functors
$\mathbf{B}M \to \mathbf{B}N$
(Proposition~\ref{ch02:prop:functor-monoid}). A natural transformation
$f \Rightarrow g$ has a single component: an element
$n \in N = \Hom(\ast, \ast)$. Naturality at $m \in M$ reads
$g(m)\, n = n\, f(m)$ --- the element $n$ \emph{intertwines} the two
homomorphisms. So natural transformations here are not automatic at all:
they exist only when a genuine algebraic condition is solvable. For
instance, with $M = N = (\mathbb{N}, +)$ and $f, g$ the homomorphisms
``multiply by $2$'' and ``multiply by $3$,'' naturality demands
$3m + n = n + 2m$ for all $m$, i.e.\ $m = 0$ for all $m$, which fails:
there is no natural transformation from $f$ to $g$. The commuting square
has teeth.
\end{example}

\begin{example}[Morphisms of presheaves on the two-element chain]\label{ch02:ex:nat-two-chain}
Anticipating Section~\ref{ch02:sec:presheaves}: a contravariant
$\Set$-valued functor on $[2]$ amounts to a function
$r : S_1 \to S_0$ (Section~\ref{ch02:sec:presheaves} makes this precise),
and a natural transformation between two such, $r : S_1 \to S_0$ and
$r' : S_1' \to S_0'$, is a pair of functions
$\alpha_1 : S_1 \to S_1'$, $\alpha_0 : S_0 \to S_0'$ with
$r' \circ \alpha_1 = \alpha_0 \circ r$ --- a commuting square of sets.
Everything in Part~II about maps of presheaves is this picture with more
objects.
\end{example}

\section{Opposite categories and contravariance}\label{ch02:sec:opposite}

\begin{definition}[Opposite category]\label{ch02:def:opposite}
The \emph{opposite category} $\cC^{\mathrm{op}}$ of a category $\cC$ has
the same objects, and
$\Hom_{\cC^{\mathrm{op}}}(A, B) = \Hom_{\cC}(B, A)$, with composition
inherited from $\cC$ in reversed order and the same identities. Clearly
$(\cC^{\mathrm{op}})^{\mathrm{op}} = \cC$.
\end{definition}

\begin{example}\label{ch02:ex:opposite-poset}
For a poset $(P, \le)$, the opposite category is the poset
$(P, \ge)$: all arrows reverse. For the powerset lattice
$\mathcal{P}(X)$ ordered by inclusion, the opposite is $\mathcal{P}(X)$
ordered by reverse inclusion --- which is again a powerset lattice via
complementation, though we never need that fact.
\end{example}

\begin{definition}[Contravariant functor]\label{ch02:def:contravariant}
A \emph{contravariant functor} from $\cC$ to $\mathcal{D}$ is a functor
$\cC^{\mathrm{op}} \to \mathcal{D}$. Concretely: it assigns to each
object $A$ an object $F(A)$, and to each morphism $f : A \to B$ of $\cC$
a morphism $F(f) : F(B) \to F(A)$ \emph{in the reverse direction}, with
$F(g \circ f) = F(f) \circ F(g)$ and $F(\id_A) = \id_{F(A)}$.
\end{definition}

The duality principle --- every theorem about categories yields a dual
theorem by applying it to opposite categories --- is used lightly in this
monograph, mainly through the mono/epi and section/retraction dualities
already visible in Proposition~\ref{ch02:prop:split-epi-is-epi}. What we
need constantly is the contravariant direction itself: presheaves, the
subject of the next section, are contravariant by definition, and the
directional bookkeeping (restriction goes \emph{down}) must become
second nature.

\section{Presheaves}\label{ch02:sec:presheaves}

\begin{definition}[Presheaf]\label{ch02:def:presheaf}
Let $\cC$ be a small category. A \emph{presheaf} on $\cC$ is a functor
$F : \cC^{\mathrm{op}} \to \Set$. Unwinding: $F$ assigns to each object
$C$ a set $F(C)$ (the \emph{sections over} $C$) and to each morphism
$f : D \to C$ a function $F(f) : F(C) \to F(D)$ (the \emph{restriction
along} $f$), such that restriction along an identity is the identity and
restriction along a composite is the composite of the restrictions in
reverse order. The category of presheaves is
$\PSh(\cC) = [\cC^{\mathrm{op}}, \Set]$, with natural transformations as
morphisms.
\end{definition}

When $\cC$ is a poset and $U \le V$, we write the restriction of a
section $s \in F(V)$ along the unique morphism $U \to V$ as
$s|_U \in F(U)$. The functor laws then read: $s|_V = s$ for
$s \in F(V)$, and $(s|_V)|_U = s|_U$ whenever $U \le V \le W$ and
$s \in F(W)$ --- restriction is transitive and unital, exactly the
algebra of ``forgetting part of the picture.''

\begin{example}[Presheaves on the two-element chain]\label{ch02:ex:psh-two-chain}
A presheaf $F$ on $[2]$ consists of two sets $F(1)$, $F(0)$ and one
restriction function $F(1) \to F(0)$ (along $0 \le 1$); the functor laws
add nothing further, since the only composites in $[2]$ involve
identities. So $\PSh([2])$ is the category whose objects are functions
between sets and whose morphisms are commuting squares
(Example~\ref{ch02:ex:nat-two-chain}). This tiny observation generalizes: a
presheaf on a poset is a system of sets wired together
by restriction functions along the order, one function per covering
relation, with commutativity along paths for free in chains but a real
condition once the poset has diamonds.
\end{example}

\begin{example}[Presheaves on the vee]\label{ch02:ex:psh-vee}
A presheaf on $\Lambda$ ($x \le t \ge y$) is a \emph{span}: three sets
and two functions $F(t) \to F(x)$, $F(t) \to F(y)$ with common domain.
There are no composable nonidentity pairs, hence no commutativity
conditions.
\end{example}

\begin{example}[Presheaves on the square]\label{ch02:ex:psh-square}
A presheaf on $B = \mathcal{P}(\{1,2\})$ consists of four sets
$F(\{1,2\})$, $F(\{1\})$, $F(\{2\})$, $F(\varnothing)$ and five
restriction functions (one per nonidentity inclusion), subject to
\emph{one} nontrivial commutativity condition: the two paths from
$F(\{1,2\})$ to $F(\varnothing)$ --- via $F(\{1\})$ and via $F(\{2\})$
--- must agree, because both equal restriction along
$\varnothing \subseteq \{1,2\}$. This is
Proposition~\ref{ch02:prop:poset-commute} biting for the first time: in
the source poset the two paths are the same morphism, so the presheaf
must send them to the same function.
\end{example}

\begin{example}[The bit presheaf]\label{ch02:ex:bit-presheaf}
On $B$, define $G(V) = \{0,1\}^V$, the set of all functions
$V \to \{0,1\}$, with restriction $s|_U$ the literal restriction of the
function $s$ to the subset $U$. Then $G(\varnothing)$ is a singleton
(the empty function), $|G(\{1\})| = |G(\{2\})| = 2$, and
$|G(\{1,2\})| = 4$. Functoriality is immediate: restricting a function
to $V$ then to $U \subseteq V$ is restricting it to $U$. This presheaf
--- ``a bit written on each point in view'' --- is the germ of the
witness presheaves of Chapter~9, and the reader met it informally in the
vignette of Chapter~1.
\end{example}

\begin{example}[The constant-verdict presheaf]\label{ch02:ex:verdict-presheaf}
On $B$, define $H(V) = \{\mathsf{p}, \mathsf{q}\}$ for all nonempty $V$
and $H(\varnothing) = \{\ast\}$, with identity restrictions among
nonempty contexts and the unique map to $H(\varnothing)$. Functoriality
is a short check (all composites of identities and unique-to-singleton
maps agree). This is the second presheaf of the Chapter~1 vignette; its
failure to be a sheaf for the union cover of $\{1,2\}$ is worked in
Chapter~4's terms once the sheaf condition is available.
\end{example}

\begin{example}[Constant, terminal, and initial presheaves]\label{ch02:ex:constant-presheaf}
For any set $S$, the \emph{constant presheaf} $\Delta S$ assigns
$S$ to every object and $\id_S$ to every morphism. Two special cases
recur. The \emph{terminal presheaf} $\mathbf{1} = \Delta\{\ast\}$ has
exactly one section everywhere; every presheaf admits exactly one
morphism to it (each component must be the unique map to a
singleton), so it is terminal in $\PSh(\cC)$. The \emph{initial
presheaf} $\Delta\varnothing$ has no sections anywhere and admits
exactly one morphism to every presheaf (all components empty
functions). Beware a poset-specific subtlety used with force in
Chapter~10: the terminal presheaf on $B$ has one section over
$\varnothing$, but a general presheaf may have zero, two, or fifty
sections over the bottom context --- $F(\varnothing)$ is data, not
convention, and constructions that need a singleton there (as the
explicit reflector of Chapter~10 does) must impose it explicitly
rather than assume it.
\end{example}

\begin{example}[Presheaves on a discrete poset]\label{ch02:ex:psh-discrete}
Let $P$ be the discrete two-element poset (two objects, no relation
besides reflexivity). A presheaf on $P$ is just a pair of sets with
no connecting data, and a morphism is a pair of functions:
$\PSh(P) \cong \Set \times \Set$. Discreteness kills all interaction
--- there are no restrictions to be compatible with. This degenerate
case is the shape of the
\emph{resolved} world of Part~II: after the joined topology
discretizes the context factor, the resolved subcategory behaves like
a product of independent copies, one per atom, with no remaining
interaction --- the expressivity boundary discussed in Chapters~1
and~12. Interaction lives in the order; flatten the order and only
juxtaposition remains.
\end{example}

\begin{exercise}\label{ch02:exr:count-presheaves}
Count the presheaves on $[2]$ whose section sets are subsets of a fixed
two-element set $\{a, b\}$. (There are four choices for each of
$F(1), F(0)$ from $\{\varnothing, \{a\}, \{b\}, \{a,b\}\}$ and then
$|F(0)|^{|F(1)|}$ choices of restriction; summing over the sixteen
pairs, with the convention that there is exactly one function
$\varnothing \to \varnothing$, gives $4 + 2 \cdot 4 + 6 = 18$.
Carrying the sum out by hand --- rather than reading the total off
--- is a useful five-minute drill in what a presheaf \emph{is}.)
\end{exercise}

\section{Representable presheaves}\label{ch02:sec:representables}

\begin{definition}[Representable presheaf]\label{ch02:def:representable}
For an object $C$ of a small category $\cC$, the \emph{representable
presheaf} $yC$ is
\[
yC = \Hom_{\cC}(-, C) : \cC^{\mathrm{op}} \to \Set,
\]
sending $D \mapsto \Hom_{\cC}(D, C)$ and a morphism $h : D' \to D$ to
the precomposition map $f \mapsto f \circ h$. A presheaf is
\emph{representable} if it is naturally isomorphic to some $yC$. The
assignment $C \mapsto yC$ extends to a functor
$y : \cC \to \PSh(\cC)$ (the \emph{Yoneda embedding}): a morphism
$g : C \to C'$ induces the natural transformation
$yg : yC \Rightarrow yC'$ with components ``postcompose with $g$.''
\end{definition}

Functoriality of $yC$ is associativity of composition; naturality of
$yg$ is associativity again. The representable $yC$ is the category's
own record of ``all ways of mapping into $C$'' --- and sieves, the
central objects of Chapter~3, will be exactly the subfunctors of
representables.

\begin{proposition}[Representables on a poset]\label{ch02:prop:representable-poset}
Let $P$ be a poset and $W \in P$. Then for every $V \in P$,
\[
yW(V) =
\begin{cases}
\{\ast\} & \text{if } V \le W,\\
\varnothing & \text{otherwise,}
\end{cases}
\]
and all restriction maps are the unique possible functions. Thus $yW$ is
the ``indicator presheaf'' of the principal down-set
$\mathord{\downarrow} W = \{V : V \le W\}$: it has exactly one section
over each context below $W$ and none elsewhere.
\end{proposition}

\begin{proof}
$yW(V) = \Hom_P(V, W)$, which is a singleton if $V \le W$ and empty
otherwise, by definition of the poset category. Restriction along
$U \le V$ maps $\Hom_P(V, W) \to \Hom_P(U, W)$; when both are singletons
it is the unique function between singletons, when the source is empty
it is the empty function, and the case source-nonempty,
target-empty cannot occur since $U \le V \le W$ implies $U \le W$.
Functoriality holds because all diagrams of functions between sets with
at most one element commute.
\end{proof}

\begin{example}\label{ch02:ex:representables-square}
On $B = \mathcal{P}(\{1,2\})$, the four representables have section
counts (over $\varnothing, \{1\}, \{2\}, \{1,2\}$ in that order):
$y\varnothing : (1, 0, 0, 0)$; $y\{1\} : (1, 1, 0, 0)$;
$y\{2\} : (1, 0, 1, 0)$; $y\{1,2\} : (1, 1, 1, 1)$. Each is the
indicator of a down-set, as
Proposition~\ref{ch02:prop:representable-poset} says. Note $y\{1,2\}$ is
the \emph{terminal} presheaf on $B$ (one section everywhere) --- in
general, on a poset with a top element, the representable of the top is
terminal.
\end{example}

\section{Subpresheaves}\label{ch02:sec:subpresheaves}

Chapter~3's central definition --- the sieve --- will be introduced twice:
once elementarily, as a precomposition-closed set of morphisms, and once
structurally, as a subfunctor of a representable. The structural half needs
the notion of subpresheaf, which we set up now and practice on the small
posets.

\begin{definition}[Subpresheaf]\label{ch02:def:subpresheaf}
Let $F$ be a presheaf on a small category $\cC$. A \emph{subpresheaf}
(or \emph{subfunctor}) $S \subseteq F$ is a choice of subset
$S(C) \subseteq F(C)$ for each object $C$ such that restrictions
preserve membership: for every morphism $f : D \to C$ and every
$s \in S(C)$, one has $F(f)(s) \in S(D)$. Equivalently, $S$ is itself a
presheaf (with the restricted restriction maps) and the inclusions
$S(C) \subseteq F(C)$ form a natural transformation with injective
components.
\end{definition}

The closure condition is the whole content: a subpresheaf is not an
arbitrary family of subsets but a family stable under restriction ---
``once a section is in, all its shadows are in.'' On a poset the
condition reads: if $s \in S(V)$ and $U \le V$ then $s|_U \in S(U)$.

\begin{example}[Subpresheaves of a representable on the vee]\label{ch02:ex:subrep-vee}
Work on $\Lambda$ ($x, y \le t$) and take $F = yt$, which has exactly
one section over each of $x$, $y$, $t$
(Section~\ref{ch02:sec:lab-vee}). A subpresheaf $S \subseteq yt$
chooses, for each object, either the empty set or the full singleton,
subject to restriction-stability. If $S(t) = \{\ast\}$, then
restricting the section of $yt(t)$ along $x \le t$ and $y \le t$
forces $S(x) = S(y) = \{\ast\}$: the only option is $S = yt$ itself.
If $S(t) = \varnothing$, the choices at $x$ and $y$ are free and
independent, since no nonidentity morphism has domain above $x$ or
$y$. That gives four subpresheaves:
$(\varnothing, \varnothing, \varnothing)$,
$(\{\ast\}, \varnothing, \varnothing)$,
$(\varnothing, \{\ast\}, \varnothing)$,
$(\{\ast\}, \{\ast\}, \varnothing)$, writing values over
$(x, y, t)$. Total: five subpresheaves of $yt$. Chapter~3 will
recognize these five as exactly the five \emph{sieves} on the object
$t$ of $\Lambda$, and the forcing argument just made --- ``a full top
forces full shadows'' --- as the elementary statement that a sieve
containing the identity is maximal.
\end{example}

\begin{example}[Subpresheaves of the bit presheaf]\label{ch02:ex:sub-bit}
On $B$, let $G$ be the bit presheaf
(Example~\ref{ch02:ex:bit-presheaf}) and define
$S(V) \subseteq G(V)$ to consist of the \emph{constant} functions
$V \to \{0,1\}$ (and the empty function over $\varnothing$). Then $S$
is a subpresheaf: a restriction of a constant function is constant.
Section counts: $(1, 2, 2, 2)$ over
$(\varnothing, \{1\}, \{2\}, \{1,2\})$. By contrast, the family
$T(V) = \{$functions summing to $1$ over $V\}$ (interpreting the sum
of the empty function as $0$) is \emph{not} a subpresheaf: the
function on $\{1,2\}$ with values $1, 0$ sums to $1$, but its
restriction to $\{2\}$ is the zero function, which sums to $0$ and so
fails membership. Checking restriction-stability is a real check, and
Chapter~11's certificate performs precisely this kind of check ---
mechanically, for every candidate --- under the name functoriality
validation.
\end{example}

\begin{exercise}\label{ch02:exr:sub-two-chain}
Enumerate all subpresheaves of $y1$ on the chain $[2]$. (There are
three; compare them with the three sieves on the object $1$ found in
Chapter~3.)
\end{exercise}

\section{The Yoneda lemma}\label{ch02:sec:yoneda}

The Yoneda lemma says that mapping out of a representable is the same as
evaluating: natural transformations $yC \Rightarrow F$ correspond to
elements of $F(C)$. It is the fundamental bookkeeping theorem of
presheaf categories, and in Part~II it is used chiefly in two ways:
through the identification of sieves with subfunctors of representables
(Chapter~3), and through the sheaf condition's formulation as a
bijection $F(C) \cong \mathrm{Nat}(S, F)$ for covering sieves $S$
(Chapter~4), which is a Yoneda-style statement with the sieve in place
of the full representable.

\begin{theorem}[Yoneda lemma; statement]\label{ch02:thm:yoneda}
Let $\cC$ be a small category, $C \in \Ob(\cC)$, and
$F \in \PSh(\cC)$. The map
\[
\Phi_{C,F} : \mathrm{Nat}(yC, F) \longrightarrow F(C),
\qquad \alpha \longmapsto \alpha_C(\id_C),
\]
is a bijection, natural in both $C$ and $F$. In particular
(taking $F = yC'$) the Yoneda embedding $y : \cC \to \PSh(\cC)$ is
fully faithful. See \cite{riehl,awodey,maclanecwm} for the general
proof.
\end{theorem}

We prove the poset case in complete detail; it is the case Part~II
uses, and its proof displays the mechanism --- a natural family is
determined by its value at the ``generic'' point --- without the
notational overhead of the general case.

\begin{theorem}[Yoneda lemma for posets; full proof]\label{ch02:thm:yoneda-poset}
Let $P$ be a poset, $W \in P$, and $F \in \PSh(P)$. Then there is a
bijection
\[
\mathrm{Nat}(yW, F) \;\cong\; F(W),
\]
given by evaluating at $W$; it is natural in $W$ and in $F$.
\end{theorem}

\begin{proof}
Recall (Proposition~\ref{ch02:prop:representable-poset}) that $yW(V)$
is a singleton $\{\ast\}$ for $V \le W$ and empty otherwise.

\emph{The map.} Given a natural transformation
$\alpha : yW \Rightarrow F$, its component at $W$ is a function
$\alpha_W : yW(W) \to F(W)$ with $yW(W) = \{\ast\}$ (as $W \le W$);
define $\Phi(\alpha) = \alpha_W(\ast) \in F(W)$.

\emph{$\Phi$ is injective.} Suppose $\Phi(\alpha) = \Phi(\beta)$; we
show $\alpha = \beta$ componentwise. At $V \not\le W$ both components
are the unique (empty) function $\varnothing \to F(V)$. At $V \le W$,
apply naturality of $\alpha$ to the morphism $V \to W$: the square
\[
F(V \le W)\bigl(\alpha_W(\ast)\bigr) = \alpha_V\bigl(\ast|_V\bigr)
= \alpha_V(\ast)
\]
(the restriction map of $yW$ between the singletons $yW(W)$ and
$yW(V)$ sends $\ast$ to $\ast$). Hence
$\alpha_V(\ast) = \Phi(\alpha)|_V$ is determined by $\Phi(\alpha)$;
the same computation for $\beta$ gives
$\beta_V(\ast) = \Phi(\beta)|_V$. Since $\Phi(\alpha) = \Phi(\beta)$,
all components agree, so $\alpha = \beta$.

\emph{$\Phi$ is surjective.} Given $s \in F(W)$, define a candidate
transformation $\alpha^s$ by
\[
\alpha^s_V(\ast) = s|_V \quad \text{for } V \le W,
\]
and the empty function at $V \not\le W$. Naturality must be checked at
each morphism $U \le V$ of $P$. If $V \le W$ (hence $U \le W$): going
one way around the square, $\ast \in yW(V)$ maps to
$\alpha^s_V(\ast) = s|_V$ and then restricts to $(s|_V)|_U = s|_U$
(functoriality of $F$); going the other way, $\ast$ restricts to
$\ast \in yW(U)$ and then maps to $\alpha^s_U(\ast) = s|_U$. Equal. If
$V \not\le W$, the source $yW(V)$ is empty and the square commutes
vacuously. So $\alpha^s$ is natural, and
$\Phi(\alpha^s) = s|_W = s$.

\emph{Naturality in $F$.} Let $\tau : F \Rightarrow G$ be a morphism
of presheaves. We must show that for every
$\alpha \in \mathrm{Nat}(yW, F)$,
$\Phi(\tau \circ \alpha) = \tau_W(\Phi(\alpha))$; and indeed
$(\tau \circ \alpha)_W(\ast) = \tau_W(\alpha_W(\ast))$ by definition
of composition of natural transformations.

\emph{Naturality in $W$.} Let $W \le W'$, inducing
$yW \Rightarrow yW'$ and, by precomposition, a map
$\mathrm{Nat}(yW', F) \to \mathrm{Nat}(yW, F)$. Chasing the point:
$\alpha' \in \mathrm{Nat}(yW', F)$ precomposes to the transformation
whose value at $W$ on $\ast$ is $\alpha'_W(\ast)$, and the injectivity
computation above (applied to $\alpha'$ and the morphism $W \le W'$)
gives $\alpha'_W(\ast) = \alpha'_{W'}(\ast)|_W = \Phi(\alpha')|_W$.
So the square comparing $\mathrm{Nat}(y(-), F)$ with $F(-)$ along
$W \le W'$ commutes: both routes send $\alpha'$ to
$\Phi(\alpha')|_W$.
\end{proof}

\begin{corollary}[Yoneda embedding on a poset]\label{ch02:cor:yoneda-poset}
For a poset $P$ and $U, W \in P$:
\[
\mathrm{Nat}(yU, yW) \;\cong\; yW(U) \;=\;
\begin{cases}
\{\ast\} & U \le W,\\
\varnothing & \text{otherwise.}
\end{cases}
\]
Hence the Yoneda embedding $y : P \to \PSh(P)$ is fully faithful ---
it is an order-embedding of $P$ into the presheaf category --- and
$yU \cong yW$ implies $U = W$.
\end{corollary}

\begin{proof}
The displayed bijection is Theorem~\ref{ch02:thm:yoneda-poset} with
$F = yW$. Full faithfulness is the statement that
$\Hom_P(U, W) \to \mathrm{Nat}(yU, yW)$ is a bijection, and both sides
are singletons or both empty according to whether $U \le W$. If
$yU \cong yW$, then in particular $\mathrm{Nat}(yU, yW)$ and
$\mathrm{Nat}(yW, yU)$ are both nonempty, so $U \le W$ and $W \le U$,
whence $U = W$ by antisymmetry.
\end{proof}

\begin{example}[Yoneda verified by hand on the two-element
chain]\label{ch02:ex:yoneda-hand}
Let $F$ be any presheaf on $[2]$, i.e.\ a function
$r : F(1) \to F(0)$. Compute $\mathrm{Nat}(y0, F)$ directly: the
representable $y0$ has $y0(0) = \{\ast\}$ and $y0(1) = \varnothing$. A
natural transformation $y0 \Rightarrow F$ is a choice of
$\alpha_0(\ast) \in F(0)$ and the empty function at $1$; the only
naturality square (along $0 \le 1$) has empty upper-left corner, hence
commutes. So $\mathrm{Nat}(y0, F) \cong F(0)$: one transformation per
element of $F(0)$, as the lemma predicts. Now $\mathrm{Nat}(y1, F)$:
$y1$ has a single section over each object. A transformation is a pair
$\alpha_1(\ast) = s \in F(1)$, $\alpha_0(\ast) = u \in F(0)$, and the
naturality square along $0 \le 1$ forces $u = r(s)$. So the data
reduce to the free choice of $s \in F(1)$:
$\mathrm{Nat}(y1, F) \cong F(1)$. Again as predicted --- and the
mechanism is visible: the component at the representing object is
free, and naturality propagates it downward.
\end{example}

\begin{exercise}\label{ch02:exr:yoneda-square}
On $B$, take $F = G$ the bit presheaf of
Example~\ref{ch02:ex:bit-presheaf}. Verify by direct enumeration that
$\mathrm{Nat}(y\{1\}, G)$ has exactly two elements and that they
correspond to the two elements of $G(\{1\}) = \{0,1\}^{\{1\}}$ as in
Theorem~\ref{ch02:thm:yoneda-poset}.
\end{exercise}

\begin{example}[Counting all morphisms between two presheaves]\label{ch02:ex:count-nat}
The Yoneda lemma computes $\mathrm{Nat}(yW, F)$ in one stroke; for
non-representable sources one counts squares directly, and it is
instructive to do so once. On $[2]$, let $F$ be the presheaf with
$F(1) = \{s, s'\}$, $F(0) = \{u\}$ (restriction the constant map)
and let $F'$ be the presheaf with $F'(1) = \{v\}$,
$F'(0) = \{w, w'\}$, restriction $v \mapsto w$. A morphism
$\alpha : F \Rightarrow F'$ is a pair of functions
$\alpha_1 : \{s, s'\} \to \{v\}$ and
$\alpha_0 : \{u\} \to \{w, w'\}$ making the square commute:
$\alpha_0(u)$ must equal the restriction of $\alpha_1(s)$, which is
$w$. So $\alpha_1$ is forced (constant at $v$), $\alpha_0$ is forced
($u \mapsto w$), and $\mathrm{Nat}(F, F') $ has exactly one element
--- even though there are $2$ functions $\{u\} \to \{w, w'\}$, the
square eliminates one. In the other direction,
$\mathrm{Nat}(F', F)$: $\alpha_1 : \{v\} \to \{s, s'\}$ is a free
binary choice, and $\alpha_0 : \{w, w'\} \to \{u\}$ is forced; the
square commutes either way (both routes land in the singleton
$\{u\}$), so there are exactly $2$ morphisms. Naturality is a
genuine constraint whose bite depends on direction --- a five-minute
computation worth repeating on any pair of small presheaves the
reader invents.
\end{example}

\begin{remark}[Covariant representables]\label{ch02:rem:covariant-yoneda}
Everything in this section dualizes: the covariant functor
$\Hom_{\cC}(C, -)$ is a ``co-presheaf,'' and the dual Yoneda lemma
identifies $\mathrm{Nat}(\Hom_{\cC}(C, -), F) \cong F(C)$ for
covariant $F$. This monograph works exclusively with the
contravariant side --- evidence restricts downward --- and we will
not use the dual statement, but the reader consulting
\cite{maclanecwm} or \cite{riehl} should be aware that both
conventions circulate and that translation is mechanical.
\end{remark}

\begin{remark}[On the general proof]\label{ch02:rem:yoneda-general}
The general Yoneda proof is the same two-step dance with $\id_C$
playing the role of the generic point $\ast$: injectivity because
naturality forces $\alpha_D(f) = F(f)(\alpha_C(\id_C))$ for every
$f : D \to C$, surjectivity because that formula, read as a
definition, is natural. The poset proof above \emph{is} the general
proof specialized; nothing was simplified away except the need to
track which morphism one is restricting along, since in a poset there
is never more than one. See \cite[Ch.~2]{riehl} or
\cite[Ch.~8]{awodey}.
\end{remark}

\section{Limits and pullbacks}\label{ch02:sec:limits}

Limits formalize ``the universal object mapping compatibly into a
diagram.'' Part~II needs three instances --- terminal objects, products,
and above all pullbacks --- plus the fact that in a poset all of them
reduce to order-theoretic infima. We give the general definitions and
then specialize.

\begin{definition}[Diagram, cone, limit]\label{ch02:def:limit}
Let $\mathcal{J}$ be a small category (the \emph{shape}) and
$D : \mathcal{J} \to \cC$ a functor (a \emph{diagram} of shape
$\mathcal{J}$ in $\cC$). A \emph{cone} over $D$ with vertex
$Z \in \Ob(\cC)$ is a family of morphisms
$(p_j : Z \to D(j))_{j \in \Ob(\mathcal{J})}$ such that for every
morphism $u : j \to j'$ of $\mathcal{J}$,
$D(u) \circ p_j = p_{j'}$. A \emph{limit} of $D$ is a cone
$(\lim D, (\pi_j))$ with the universal property: for every cone
$(Z, (p_j))$ there is a unique morphism $u : Z \to \lim D$ with
$\pi_j \circ u = p_j$ for all $j$.
\end{definition}

\begin{proposition}[Limits are unique up to unique isomorphism]\label{ch02:prop:limit-unique}
Any two limits of the same diagram are isomorphic by a unique
isomorphism commuting with the cone maps.
\end{proposition}

\begin{proof}
Let $(P, (\pi_j))$ and $(P', (\pi'_j))$ both be limits. The cone
$(P', (\pi'_j))$ factors through $(P, (\pi_j))$ by a unique
$u : P' \to P$, and symmetrically $(P, (\pi_j))$ factors by a unique
$v : P \to P'$. The composite $u \circ v : P \to P$ is a factorization
of the cone $(P, (\pi_j))$ through itself; but $\id_P$ is another such
factorization, and uniqueness in the universal property forces
$u \circ v = \id_P$. Symmetrically $v \circ u = \id_{P'}$.
\end{proof}

\begin{example}[Terminal object, product]\label{ch02:ex:terminal-product}
The limit of the empty diagram is a \emph{terminal object} $1$: an
object with exactly one morphism from every object. In $\Set$, any
singleton is terminal. The limit of a discrete two-object diagram
$(A, B)$ is the \emph{product} $A \times B$ with its projections; in
$\Set$, the cartesian product. In a poset, a terminal object is a top
element, and a product of $U, V$ is a \emph{meet} (greatest lower
bound) $U \wedge V$: the universal property of the product says
exactly ``$Z \le U$ and $Z \le V$ iff $Z \le U \wedge V$.''
\end{example}

\begin{definition}[Pullback]\label{ch02:def:pullback}
A \emph{cospan} in $\cC$ is a diagram
$A \xrightarrow{\,f\,} W \xleftarrow{\,g\,} B$. A \emph{pullback} of
the cospan is a limit of it: an object $P$ with morphisms
$p : P \to A$, $q : P \to B$ satisfying $f \circ p = g \circ q$, such
that for any $Z$ with $a : Z \to A$, $b : Z \to B$ and
$f \circ a = g \circ b$ there is a unique $u : Z \to P$ with
$p \circ u = a$ and $q \circ u = b$. (The cone leg to $W$ is
determined as $f \circ p$, so one omits it.) Notation:
$P = A \times_W B$.
\end{definition}

\begin{proposition}[Pullbacks in $\Set$]\label{ch02:prop:pullback-set}
In $\Set$, every cospan $A \xrightarrow{f} W \xleftarrow{g} B$ has a
pullback:
\[
A \times_W B = \{(a, b) \in A \times B \mid f(a) = g(b)\},
\]
with the coordinate projections.
\end{proposition}

\begin{proof}
The projections satisfy $f \circ p = g \circ q$ by definition of the
set. Given a competing cone $(Z, a, b)$ with $f \circ a = g \circ b$,
define $u : Z \to A \times_W B$ by $u(z) = (a(z), b(z))$; this lands
in the pullback set precisely because $f(a(z)) = g(b(z))$, and it is
the unique function whose postcompositions with the projections are
$a$ and $b$, since a function into a subset of a product is determined
by its coordinates.
\end{proof}

\begin{example}[Pullback of subset inclusions is intersection]\label{ch02:ex:pullback-subsets}
Let $A, B \subseteq W$ be subsets with inclusion maps
$A \hookrightarrow W \hookleftarrow B$. The pullback set of
Proposition~\ref{ch02:prop:pullback-set} is
$\{(a, b) \mid a = b \in A \cap B\}$, which the map
$(a,b) \mapsto a$ bijects with $A \cap B$; under this bijection the
projections become the inclusions of $A \cap B$ into $A$ and $B$. So
in $\Set$, ``pullback of monos'' computes intersection --- the
prototype of the next proposition, which is the version Part~II
actually uses.
\end{example}

\begin{proposition}[Pullbacks in a poset are meets]\label{ch02:prop:pullback-poset}
Let $P$ be a poset and $U \to W \leftarrow V$ a cospan in $P$ (i.e.\
$U \le W$ and $V \le W$). A pullback of this cospan exists if and only
if the meet $U \wedge V$ exists in $P$, and in that case the pullback
is $U \wedge V$ with the unique morphisms to $U$ and $V$.
\end{proposition}

\begin{proof}
A cone over the cospan with vertex $Z$ is a pair of morphisms
$Z \to U$, $Z \to V$ whose composites to $W$ agree; the agreement is
automatic (Proposition~\ref{ch02:prop:poset-commute}), so a cone is
exactly an element $Z$ with $Z \le U$ and $Z \le V$ --- a lower bound
of $\{U, V\}$. The universal property of the pullback says: $P^*$ is a
lower bound such that every lower bound $Z$ admits a (unique,
automatic) morphism $Z \to P^*$, i.e.\ every lower bound satisfies
$Z \le P^*$. That is precisely the definition of the greatest lower
bound. So pullbacks of the cospan and meets of $\{U, V\}$ are the same
data. Uniqueness of morphisms is automatic throughout, again by
Proposition~\ref{ch02:prop:poset-commute}.
\end{proof}

\begin{corollary}[Pullbacks in a powerset lattice are intersections]\label{ch02:cor:pullback-powerset}
In $\mathcal{P}(X)$ ordered by inclusion, every cospan
$U \subseteq W \supseteq V$ has a pullback, namely $U \cap V$.
\end{corollary}

\begin{proof}
$\mathcal{P}(X)$ has all binary meets, and the meet of $U$ and $V$ is
$U \cap V$: it is a lower bound ($U \cap V \subseteq U$,
$U \cap V \subseteq V$), and any $Z$ with $Z \subseteq U$ and
$Z \subseteq V$ satisfies $Z \subseteq U \cap V$. Apply
Proposition~\ref{ch02:prop:pullback-poset}.
\end{proof}

This corollary is used on nearly every page of Part~II: the context
lattice $L = \mathcal{P}(\Omega)$ has all pullbacks, they are
intersections, and the pullback of a covering piece $V \cap B \to V$
along $U \subseteq V$ is $(V \cap B) \cap U = U \cap B$ --- the
one-line computation at the heart of the coverage-stability lemma of
Chapter~9.

\begin{example}[A poset without a pullback]\label{ch02:ex:no-pullback}
In the vee $\Lambda$ ($x \le t \ge y$), the cospan
$x \to t \leftarrow y$ has no pullback: a pullback would be a meet
$x \wedge y$, i.e.\ a greatest lower bound of $\{x, y\}$, and
$\{x, y\}$ has \emph{no} lower bound at all in $\Lambda$ (neither $x$
nor $y$ is below the other, and $t$ is above both, not below).
Existence of pullbacks is a genuine property of a poset, not a
formality --- powerset lattices have it; arbitrary posets need not.
The registered legacy category is required to have pullback closure by
its axioms; the context lattice has pullbacks by
Corollary~\ref{ch02:cor:pullback-powerset}; the vee shows the
requirement is not vacuous.
\end{example}

\begin{exercise}\label{ch02:exr:pullback-square}
In $B = \mathcal{P}(\{1,2\})$, compute the pullback of the cospan
$\{1\} \to \{1,2\} \leftarrow \{2\}$ and verify explicitly that every
lower bound of $\{\{1\}, \{2\}\}$ factors through it. (Answer:
$\varnothing$; the only lower bound is $\varnothing$ itself.)
\end{exercise}

\subsection{Equalizers and general limits}\label{ch02:sec:equalizers}

For completeness we record the remaining finite-limit shape and the
standard reduction.

\begin{definition}[Equalizer]\label{ch02:def:equalizer}
Given parallel morphisms $f, g : A \to B$, an \emph{equalizer} is a
morphism $e : E \to A$ with $f \circ e = g \circ e$, universal among
such: any $h : Z \to A$ with $f \circ h = g \circ h$ factors uniquely
through $e$. In $\Set$, $E = \{a \in A \mid f(a) = g(a)\}$ with the
inclusion. In a poset, any parallel pair is equal
(Proposition~\ref{ch02:prop:poset-commute}), so every equalizer is an
identity --- equalizers are invisible in the poset world, which is
one more simplification Part~II enjoys.
\end{definition}

\begin{remark}[Limits from products and equalizers]\label{ch02:rem:limits-reduction}
A category with all small products and equalizers has all small
limits, by the standard construction of a limit as an equalizer of two
maps between products \cite[Ch.~V]{maclanecwm}; $\Set$ therefore has
all small limits. We do not use the general construction --- every
limit needed in this monograph is a terminal object, a binary product,
or a pullback, built by hand --- but it is the reason one may speak of
``the'' limit apparatus of $\Set$ without further inventory.
\end{remark}

\subsection{Limits of presheaves are computed pointwise}\label{ch02:sec:pointwise}

One structural fact about presheaf categories is needed in Chapter~10,
where an explicit reflector is built out of products of section sets:
limits of presheaves are computed objectwise.

\begin{proposition}[Pointwise limits]\label{ch02:prop:pointwise}
Let $\cC$ be small and let $D : \mathcal{J} \to \PSh(\cC)$ be a finite
(or small) diagram of presheaves. Define, for each object $C$ of
$\cC$, $(\lim D)(C) = \lim_j \bigl(D(j)(C)\bigr)$, the limit in
$\Set$ of the evaluated diagram, with restriction maps induced by the
universal property. Then $\lim D$, so defined, is a presheaf and is a
limit of $D$ in $\PSh(\cC)$.
\end{proposition}

\begin{proof}[Proof for binary products, which is the case used]
Let $F, G$ be presheaves and define
$(F \times G)(C) = F(C) \times G(C)$ with restriction
$(s, t)|_f = (F(f)(s), G(f)(t))$ for $f : D \to C$. Functoriality is
componentwise. The projections $\pi_F, \pi_G$ are natural
transformations (their naturality squares are the definitions of the
restriction just given). Given a presheaf $Z$ with natural
transformations $\alpha : Z \Rightarrow F$, $\beta : Z \Rightarrow G$,
define $u_C(z) = (\alpha_C(z), \beta_C(z))$; each $u_C$ is the unique
function through which $\alpha_C, \beta_C$ factor
(Proposition~\ref{ch02:prop:pullback-set}'s product analogue in
$\Set$), and the family $(u_C)$ is natural because its components are
built from natural components. Uniqueness of $u$ as a natural
transformation follows from uniqueness of each component. The same
argument, with the $\Set$-limit of the evaluated diagram replacing the
cartesian product, proves the general statement; see
\cite[Ch.~V]{maclanecwm} or \cite[Ch.~3]{riehl}.
\end{proof}

\begin{example}\label{ch02:ex:pointwise-square}
On $B$, the product of the bit presheaf $G$
(counts $(1,2,2,4)$) and the constant-verdict presheaf $H$ (counts
$(1,2,2,2)$) has section counts $(1, 4, 4, 8)$, each restriction
acting coordinatewise. No new phenomena: the product presheaf is just
both records kept side by side --- which is exactly how Chapter~10's
product-over-blocks reflector should be read, one factor per
partition block.
\end{example}

\subsection{Equality versus natural isomorphism}\label{ch02:sec:eq-vs-iso}

A final calibration, load-bearing for Chapter~10. Two functors can be
naturally isomorphic without being equal, and the discipline of this
monograph keeps the distinction visible. A representative example: on
$B$, the presheaf $G \times H$ of
Example~\ref{ch02:ex:pointwise-square} and the presheaf
$H \times G$ are not literally the same functor --- their sections are
ordered pairs with opposite component order --- but the swap maps
$(s, t) \mapsto (t, s)$ assemble into a natural isomorphism between
them. All structural claims in this monograph that identify two
constructed functors are claims of natural isomorphism with named
comparison maps, never claims of literal equality; the composition
theorem of Chapter~10 states this explicitly for its nested-product
flattenings, since a nested tuple $((a, b), (c, d))$ and a flat tuple
indexed by four atoms are different set-theoretic objects that a
chosen implementation must convert between. The certificate of
Chapter~11 checks the conversion maps --- both directions, inverse
laws, naturality --- rather than pretending the objects coincide.
Acquire the habit of asking ``equal, or isomorphic via which map?''
now, on examples this small, where the answer is always computable by
hand.

\section{The small-poset laboratory}\label{ch02:sec:laboratory}

This closing section consolidates the running examples into a single
reference, tabulating for each of the four standard posets the data
that Chapters~3 and~9 will call on. Everything below is a restatement
or short extension of computations already made; the point of
gathering them is that Part~II's sixteen-object arguments are these
arguments with a larger index set, and a reader who owns the small
cases owns the method.

\subsection{The two-element chain}\label{ch02:sec:lab-two-chain}

Objects $0 \le 1$; morphisms $\id_0, \id_1, f : 0 \to 1$. (For an
$n$-element chain the morphism count is the number of pairs
$i \le j$, namely $n(n+1)/2$ --- here $3$; the $3^{|X|}$ formula of
Proposition~\ref{ch02:prop:threepow} is specific to powerset
lattices.) A presheaf is a function
$F(1) \to F(0)$; a morphism of presheaves is a commuting square.
Representables: $y0 = (1, 0)$ and $y1 = (1, 1)$, writing section
counts over $(0, 1)$. Yoneda: verified by hand in
Example~\ref{ch02:ex:yoneda-hand}. Pullbacks: every cospan in a chain
has a pullback, since any two elements of a chain are comparable and
the meet is the smaller; in particular
$0 \to 1 \leftarrow 0$ pulls back to $0$.

\subsection{The three-element chain}\label{ch02:sec:lab-three-chain}

Objects $0 \le 1 \le 2$; six morphisms; the unique nontrivial
composition $(1 \to 2) \circ (0 \to 1) = (0 \to 2)$. A presheaf is a
composable pair of functions $F(2) \to F(1) \to F(0)$; the functor law
adds the requirement that the composite equal the restriction along
$0 \le 2$, which is a definition rather than a constraint (there is
nothing else for it to be). Representables, section counts over
$(0,1,2)$: $y0 = (1,0,0)$, $y1 = (1,1,0)$, $y2 = (1,1,1)$. The chain
illustrates transitivity of restriction in its simplest nontrivial
form and will reappear in Chapter~3 as the first example with a
three-layer sieve structure.

\subsection{The vee}\label{ch02:sec:lab-vee}

Objects $x, y \le t$; five morphisms; no nontrivial compositions. A
presheaf is a span $F(x) \leftarrow F(t) \to F(y)$. Representables,
section counts over $(x, y, t)$: $yx = (1, 0, 0)$, $yy = (0, 1, 0)$,
$yt = (1, 1, 1)$. No meet of $\{x, y\}$, hence no pullback of the
cospan $x \to t \leftarrow y$
(Example~\ref{ch02:ex:no-pullback}). The vee is the minimal
illustration that ``the overlap of two views'' need not exist as an
object --- in the powerset lattices of Part~II it always does, and the
contrast is what the pullback-closure axioms are about.

\subsection{The square}\label{ch02:sec:lab-square}

$B = \mathcal{P}(\{1,2\})$: four objects, nine morphisms
(Proposition~\ref{ch02:prop:threepow}), one nontrivial commutativity
condition per presheaf (Example~\ref{ch02:ex:psh-square}).
Representables tabulated in
Example~\ref{ch02:ex:representables-square}; the top representable is
terminal. All pullbacks exist and are intersections
(Corollary~\ref{ch02:cor:pullback-powerset}); the generic instance
$\{1\} \cap \{2\} = \varnothing$ is
Exercise~\ref{ch02:exr:pullback-square}. Two working presheaves live
here: the bit presheaf $G$ (Example~\ref{ch02:ex:bit-presheaf}), with
section counts $(1, 2, 2, 4)$ over
$(\varnothing, \{1\}, \{2\}, \{1,2\})$, and the constant-verdict
presheaf $H$ (Example~\ref{ch02:ex:verdict-presheaf}), with counts
$(1, 2, 2, 2)$. Chapter~4 shows: for the covering family
$\{ \{1\}, \{2\} \}$ of $\{1,2\}$, the presheaf $G$ glues perfectly
while $H$ does not, and the numbers $4$ versus $2 \times 2$ already
hint at why.

The square is the $|X| = 2$ member of the family
$\mathcal{P}(X)$; the context lattice $L$ of Part~II is the
$|X| = 4$ member, with $X = \Omega = \{0,1\} \times \{0,1\}$ carrying
its two coordinates. Everything tabulated here scales: $16$ objects,
$81$ morphisms, representables as down-set indicators, pullbacks as
intersections, and presheaves as commuting cubes of restriction
functions. The scaling adds bookkeeping, not ideas --- which is
precisely why the finite certificate of Chapter~11 can check the
scaled statements exhaustively, and why this chapter has spent its
pages on the small cases where the ideas are visible.

\subsection{Comparison table}\label{ch02:sec:lab-table}

\begin{center}
\begin{tabular}{lcccc}
\toprule
 & $[2]$ & $[3]$ & $\Lambda$ & $B = \mathcal{P}(\{1,2\})$ \\
\midrule
Objects & 2 & 3 & 3 & 4 \\
Morphisms & 3 & 6 & 5 & 9 \\
Top element & yes & yes & yes & yes \\
All binary meets & yes & yes & no & yes \\
All pullbacks & yes & yes & no & yes \\
Presheaf shape & function & pair & span & square \\
Commutativity conditions & 0 & 0 & 0 & 1 \\
\bottomrule
\end{tabular}
\end{center}

The context lattice $L = \mathcal{P}(\Omega)$ of Part~II extends the
last column: $16$ objects, $81$ morphisms
(Proposition~\ref{ch02:prop:threepow}), top element $\Omega$, all
meets as intersections, all pullbacks, and presheaves as commuting
assemblies of restriction functions over the inclusion order.

\subsection{Where each tool is used later}\label{ch02:sec:usage-map}

For orientation, the primer's results map onto Part~II as follows.
The poset dictionary and commutation
(Propositions~\ref{ch02:prop:poset-category}
and~\ref{ch02:prop:poset-commute}) underlie every diagram chase on
$L$. The split-epi apparatus (Section~\ref{ch02:sec:split}) is
consumed by the collapse theorem of Chapter~5; the poset immunity
remark (Remark~\ref{ch02:rem:no-split-in-posets}) explains why the
successor's site is out of its reach. Subpresheaves
(Section~\ref{ch02:sec:subpresheaves}) become sieves in Chapter~3.
Representables-as-indicators
(Proposition~\ref{ch02:prop:representable-poset}) and the poset
Yoneda lemma (Theorem~\ref{ch02:thm:yoneda-poset}) power the
subcanonicity theorem of Chapter~9 and the sheaf condition's
formulation in Chapter~4. Pullbacks-as-intersections
(Corollary~\ref{ch02:cor:pullback-powerset}) is the engine of the
coverage-stability lemma in Chapter~9. Pointwise limits
(Proposition~\ref{ch02:prop:pointwise}) legitimize the
product-over-blocks reflector of Chapter~10, and the
equality-versus-isomorphism calibration
(Section~\ref{ch02:sec:eq-vs-iso}) governs how that chapter's
composition theorem may be stated. The injections category of
Example~\ref{ch02:ex:finset} reappears as the countermodel of
Chapter~6. Nothing in the primer is idle.

\section{Summary}\label{ch02:sec:summary}

The chapter's yield, in the order Part~II consumes it: categories and
the poset dictionary (objects are contexts, morphisms are inclusions,
all diagrams commute); functors and natural transformations, with the
functor category $[\cC^{\mathrm{op}}, \Set]$ of presheaves as the
ambient world; restriction algebra on posets; representables as
down-set indicators (Proposition~\ref{ch02:prop:representable-poset});
the Yoneda lemma with full proof in the poset case
(Theorem~\ref{ch02:thm:yoneda-poset}) and its corollary that the
Yoneda embedding of a poset is an order-embedding; limits, with
pullbacks in posets identified as meets
(Proposition~\ref{ch02:prop:pullback-poset}) and in powerset lattices
as intersections (Corollary~\ref{ch02:cor:pullback-powerset}); and the
two special morphism classes --- isomorphisms (none nontrivial in
posets) and section-retraction pairs (none nontrivial in posets,
ubiquitous in $\Set$) --- whose interaction with sieves opens
Chapter~3. All of it is standard
\cite{awodey,riehl,maclanecwm,maclane}; the next chapter begins to
aim it.
